\section{Completed and Planned Evaluations}

\subsection{HotpotQA Sampled Kaggle Evaluation}

HotpotQA is now included as a sampled Kaggle evaluation, not a full benchmark claim. It is valuable because multi-hop QA stresses evidence sufficiency more directly than simple claim verification. The implementation uses cached BM25 retrieval so the 5.23M-document BEIR HotpotQA corpus is indexed once per notebook instead of once per action row.

The clean MiMo sampled runs provide the usable HotpotQA evidence in this report:

\begin{itemize}[leftmargin=*]
  \item \textbf{Full 16-action sampled matrix:} 800/800 MiMo generations succeeded on 50 sampled queries, with 80 MiMo-backed RAGAS samples. Cached BM25 retrieval gives mean recall@10 0.610 and nDCG@10 0.583. The best token-F1 row is \code{adaptive-heuristic\_\_balanced\_\_16000} at 0.097, while the best RAGAS answer-relevancy row is \code{adaptive-heuristic\_\_aggressive\_\_4000} at 0.871.
  \item \textbf{Evidence-aware budget curve:} 300/300 generations succeeded for 1k through 32k budgets, with 30 RAGAS samples. Evidence-aware context saturates after 8k in this sample; 8k is the most defensible trade-off row because it improves answer relevancy without retaining more context at 16k/32k.
  \item \textbf{Low-budget curve:} legacy and adaptive rows over 500--3000 characters completed 750/750 generations, with 75 RAGAS samples. This run tests whether the adaptive policy remains usable under aggressive context limits.
\end{itemize}

Exact match is 0.000 across the main sampled matrix, so token-F1 and judge metrics are the relevant answer-quality signals for this stage. The Groq Qwen3-32B HotpotQA run is kept as a diagnostic because provider rate limits caused 417/800 generation errors; it must not be mixed with the clean MiMo summary.

\subsection{vLLM / Vast AI Model-Bench Workflow}

The merged model-bench workflow is an operator/profiling path, not a selector-quality benchmark. It is included because deployment planning for BudgetRAG and future local Qwen experiments depends on knowing whether candidate models can be served reliably on rented GPUs. The workflow benchmarks local or existing OpenAI-compatible vLLM endpoints, records latency, time-to-first-token, tokens per second, hardware samples, and serving configuration fields such as quantization, KV-cache dtype, attention backend, speculative decoding settings, eager/CUDA-graph mode, batch limits, cache limits, and CPU offload.

The current Vast AI scripts target RTX 5060 Ti 16GB hosts with CUDA 13.0 and CUDA 12.9 profiles. Their outputs are written under ignored \code{runs/model\_bench/} directories. These measurements should be used for deployment triage and cost planning; they are not mixed with the RLAIF selector, HotpotQA, or BEIR benchmark claims.

\subsection{Retriever Diversity}

Current evidence mostly supports context-budget/action selection. To claim retrieval-context allocation, logged actions must include multiple retrieval strategies. A first retrieval-only SciFact run covers 45 jobs over 50 queries:

\begin{itemize}[leftmargin=*]
  \item 3 retrievers: BM25, graph-BM25, and hybrid-RRF;
  \item 5 policy/profile variants over 3 budgets, producing 2250 normalized action rows;
  \item every query has all 3 retrievers and all 45 action rows;
  \item generation was skipped, so this verifies logged action coverage but does not provide answer-quality or judge-label evidence.
\end{itemize}

The follow-up step was a small MiMo-\code{v2.5} generation subset with explicit \code{retriever}, \code{context\_policy}, \code{adaptive\_profile}, \code{budget}, \code{generator\_model}, and \code{judge\_model} columns. Optional web-search remains a live stress test only and must not be mixed with reproducible BEIR claims. This generation subset has now been run for 10 SciFact queries:

\begin{itemize}[leftmargin=*]
  \item 30 jobs produced 300 retriever-diverse action rows over BM25, graph-BM25, and hybrid-RRF with standard \code{mimo-v2.5}.
  \item MiMo answer judging produced 300 labels, 299 valid JSON labels, 222 labels with numeric diagnostics, and 186 clean reward rows after filtering ambiguous labels.
  \item The reward build produced 1559 preference pairs, including 1189 retrieval-context preferences.
  \item Six-seed held-out selector evaluation is only a diagnostic because the run contains 10 query groups; it verifies the pipeline but does not support a learned-selector generalization claim.
\end{itemize}

The run also exposed an operational issue: \code{MAX\_COMPLETION\_TOKENS=256} is too low for standard MiMo V2.5 generation, producing 77 empty answer strings with no request errors. Future retriever-diverse generation runs should use a larger generation cap before scaling to the full 50-query matrix.

\subsection{Larger Logged Query Groups}

The action coverage diagnostics show exact-signature sparsity. More query groups and broader logged action families are needed before a learned selector can robustly beat strong smoothed baselines.
